% !TEX program = xelatex
\documentclass[11pt,a4paper]{article}
\usepackage[margin=2.2cm,top=2cm]{geometry}
\usepackage{fontspec}
\usepackage{polyglossia}
\setdefaultlanguage{turkish}
\usepackage{amsmath,amssymb}
\usepackage{graphicx}
\usepackage{booktabs}
\usepackage{enumitem}
\usepackage{multicol}
\usepackage[most]{tcolorbox}
\usepackage{fancyhdr}
\usepackage{titlesec}
\usepackage{xcolor}
\usepackage{needspace}
\usepackage[hidelinks]{hyperref}

\definecolor{anamavi}{RGB}{0,45,98}
\definecolor{cozumyesil}{RGB}{0,110,50}
\definecolor{uyarikirmizi}{RGB}{180,20,20}

\titleformat{\section}{\large\bfseries\color{anamavi}}{\thesection.}{0.5em}{}
\titleformat{\subsection}{\normalsize\bfseries\color{anamavi}}{\thesubsection}{0.5em}{}

\pagestyle{fancy}\fancyhf{}
\lhead{\small MMAT105 Genel Matematik}
\rhead{\small Hafta 14 -- Ders Notu}
\cfoot{\small\thepage}
\renewcommand{\headrulewidth}{0.4pt}

\newtcolorbox{kural}[1][Kural]{colback=anamavi!5,colframe=anamavi,title=#1,fonttitle=\bfseries,breakable,left=4pt,right=4pt,top=2pt,bottom=2pt}
\newtcolorbox{dikkat}[1][Dikkat]{colback=uyarikirmizi!5,colframe=uyarikirmizi,title=#1,fonttitle=\bfseries,breakable,left=4pt,right=4pt,top=2pt,bottom=2pt}
\newcounter{ornek}
\newenvironment{ornek}{\refstepcounter{ornek}\par\medskip\noindent\textbf{\color{anamavi}Örnek \theornek.}\ }{\par}
\newenvironment{cozum}{\begin{tcolorbox}[colback=cozumyesil!6,colframe=cozumyesil,boxrule=0.5pt,left=4pt,right=4pt,top=2pt,bottom=2pt,breakable]\textbf{\color{cozumyesil}Çözüm.}\ }{\end{tcolorbox}}
\newcounter{alis}
\newenvironment{alistirma}{\refstepcounter{alis}\par\medskip\needspace{9\baselineskip}\noindent\textbf{\color{anamavi}Alıştırma \thealis.}\ }{\par}
\newcommand{\kendinizoncesi}{\needspace{12\baselineskip}}
\newcommand{\ab}[1]{\overline{#1}}
\setlength{\parindent}{0pt}
\setlength{\parskip}{3pt}

\begin{document}

% ---------------- Başlık ----------------
\begin{minipage}[c]{\textwidth}
{\LARGE\bfseries\color{anamavi} MMAT105 Genel Matematik}\\[4pt]
{\large Hafta 14 Ders Notu: Yüzde, Kâr-Zarar, Karışım ve\\ Faiz Problemleri, Genel Tekrar}\\[6pt]
{\small 2026--2027 Güz}
\end{minipage}

\vspace{4pt}\noindent\textcolor{anamavi}{\rule{\textwidth}{1pt}}

\textbf{Bu notun kullanımı.} Her bölümde önce kısa bir kural kutusu, ardından çözümlü örnekler yer alır. Örnekleri çözümüne bakmadan önce kendiniz deneyin. Notun sonunda önce çözümlü alıştırmalar, ardından cevap anahtarlı ``Kendiniz Çözün'' alıştırmaları ve son olarak finale hazırlık amacıyla, Hafta 1--14 konularını kapsayan karma bir tekrar bölümü yer alır. Bu notta para birimi TL'dir ve ondalık sayılar virgülle yazılır (\(12{,}5\)).

% =====================================================
\section{Yüzde}

\begin{kural}[Yüzde kavramı ve yüzde değişim]
\begin{itemize}[leftmargin=*,itemsep=1pt]
\item \(\%a=\dfrac{a}{100}\); \quad bir \(x\) sayısının \(\%a\)'sı \(x\cdot\dfrac{a}{100}\)'dür.
\item \(A\) sayısı \(B\)'nin \(\dfrac{A}{B}\cdot 100\) yüzdesidir (\(B\neq 0\)).
\item \(x\)'in \(\%a\) fazlası: \(x\left(1+\dfrac{a}{100}\right)\); \quad \(\%a\) eksiği: \(x\left(1-\dfrac{a}{100}\right)\).
\item Ardışık yüzde değişimlerde çarpanlar \textbf{çarpılır}: önce \(\%a\) artış, sonra \(\%b\) azalış \(\Rightarrow\) \(x\left(1+\dfrac{a}{100}\right)\left(1-\dfrac{b}{100}\right)\).
\end{itemize}
\end{kural}

\begin{dikkat}
Yüzdeler toplanıp çıkarılmaz: \(\%20\) artıp \(\%20\) azalan fiyat \(\%4\) azalır (\(1{,}2\cdot 0{,}8=0{,}96\)). İkinci değişim, ilk değişimden sonra oluşan \textbf{yeni} değer üzerinden hesaplanır. Sayı verilmemiş sorularda başlangıç değerine \(100\) (veya \(100x\)) demek işlemi kolaylaştırır.
\end{dikkat}

\begin{ornek}
\(450\) sayısının \(\%18\)'i kaçtır? Bulduğunuz sayı hangi sayının \(\%15\)'idir?
\end{ornek}
\begin{cozum}
\(450\cdot\dfrac{18}{100}=81\). \quad \(x\cdot\dfrac{15}{100}=81\Rightarrow x=\dfrac{8100}{15}=540\).
\end{cozum}

\begin{ornek}
\(\dfrac{a}{b}=\dfrac{3}{5}\) olduğuna göre \(a\), \(b\)'nin yüzde kaçıdır? \(a\), \(a+b\) toplamının yüzde kaçıdır?
\end{ornek}
\begin{cozum}
\(a=3k,\ b=5k\) (\(k\neq 0\)). \quad \(\dfrac{a}{b}\cdot 100=\dfrac{3}{5}\cdot 100=60\Rightarrow\%60\). \quad \(\dfrac{a}{a+b}\cdot 100=\dfrac{3}{8}\cdot 100=37{,}5\Rightarrow\%37{,}5\).
\end{cozum}

\begin{ornek}
Bir ürün önce \(\%40\) zamlanıyor, sonra yeni fiyat üzerinden \(\%25\) indirim yapılıyor. Son fiyat, ilk fiyata göre yüzde kaç değişmiştir?
\end{ornek}
\begin{cozum}
\(100\xrightarrow{+\%40}140\xrightarrow{-\%25\ (35)}105\). \ Son fiyat ilk fiyattan \(\%5\) \textbf{fazladır}. \ (Çarpanlar: \(1{,}4\cdot 0{,}75=1{,}05\).)
\end{cozum}

\begin{ornek}
Bir mağazada bir ürüne önce \(\%20\), ardından indirimli fiyat üzerinden \(\%25\) indirim yapılıyor. Bu iki indirim tek seferde yüzde kaçlık bir indirime eşittir?
\end{ornek}
\begin{cozum}
\(0{,}80\cdot 0{,}75=0{,}60\): son fiyat ilk fiyatın \(\%60\)'ı \(\Rightarrow\) toplam indirim \(\%40\) (\(\%45\) değil!).
\end{cozum}

\begin{ornek}
Bir sayının \(\%15\) fazlası ile \(\%20\) eksiği arasındaki fark \(56\)'dır. Bu sayı kaçtır?
\end{ornek}
\begin{cozum}
\(1{,}15x-0{,}80x=0{,}35x=56\Rightarrow x=160\). \ Kontrol: \(184-128=56\).
\end{cozum}

\begin{ornek}
Enflasyonun \(\%60\) olduğu bir yılda bir çalışanın maaşına \(\%20\) zam yapılıyor. Çalışanın alım gücü yüzde kaç azalmıştır?
\end{ornek}
\begin{cozum}
Başlangıçta maaş \(100\) TL, bir sepet ürün \(100\) TL olsun (maaş \(1\) sepet alıyor). Yıl sonunda maaş \(120\) TL, sepet \(160\) TL.
\[\frac{120}{160}=\frac{3}{4}=\%75\ \text{sepet}\ \Rightarrow\ \text{alım gücü } \%25\ \text{azalmıştır.}\]
\end{cozum}

\begin{ornek}
Bir dikdörtgenin uzun kenarı \(\%20\) artırılıp kısa kenarı \(\%20\) azaltılırsa alanı yüzde kaç değişir?
\end{ornek}
\begin{cozum}
Alan \(=a\cdot b\). Yeni alan: \(1{,}2a\cdot 0{,}8b=0{,}96\,ab\) \(\Rightarrow\) alan \(\%4\) \textbf{azalır}.
\end{cozum}

% =====================================================
\section{Kâr -- Zarar ve İndirim}

\begin{kural}[Kâr -- zarar]
\begin{itemize}[leftmargin=*,itemsep=1pt]
\item Satış fiyatı \(=\) Maliyet \(+\) Kâr \quad veya \quad Satış fiyatı \(=\) Maliyet \(-\) Zarar.
\item Kâr (zarar) yüzdesi, aksi söylenmedikçe \textbf{maliyet üzerinden} hesaplanır: \(\dfrac{\text{Kâr}}{\text{Maliyet}}\cdot 100\).
\item Maliyet \(M\) ise \(\%a\) kârla satış \(M\left(1+\dfrac{a}{100}\right)\), \(\%a\) zararla satış \(M\left(1-\dfrac{a}{100}\right)\).
\item İndirim yüzdesi \textbf{etiket fiyatı üzerinden} hesaplanır: \(\%d\) indirimli fiyat \(=\) Etiket \(\cdot\left(1-\dfrac{d}{100}\right)\).
\end{itemize}
\end{kural}

\begin{dikkat}
Kâr yüzdesini satış fiyatına veya etiket fiyatına bölmeyin. ``İndirimden sonra kâr edilir'' türü sorularda önce indirimli satış fiyatını bulun, sonra maliyetle karşılaştırın.
\end{dikkat}

\begin{ornek}
Bir mal, maliyetinin \(\%10\) fazlasına \(330\) TL'ye satılıyor. Maliyet kaç TL'dir?
\end{ornek}
\begin{cozum}
Maliyet \(100x\) \(\Rightarrow\) satış \(110x=330\Rightarrow x=3\). \ Maliyet \(=300\) TL.
\end{cozum}

\begin{ornek}
Bir satıcı, eşit maliyetli iki üründen birini \(\%20\) kârla, diğerini \(\%10\) zararla satıyor. Satıcı toplamda yüzde kaç kâr veya zarar etmiştir?
\end{ornek}
\begin{cozum}
Her birinin maliyeti \(100\) TL olsun. Satışlar: \(120\) ve \(90\) TL; toplam satış \(210\), toplam maliyet \(200\).
\[\frac{210-200}{200}\cdot 100=\%5\ \text{kâr}\]
\end{cozum}

\begin{ornek}
Bir ürün \(\%20\) kârla satılsaydı, \(\%10\) zararla satıldığı duruma göre \(120\) TL daha fazla para elde edilirdi. Ürünün maliyeti kaç TL'dir?
\end{ornek}
\begin{cozum}
\(1{,}2M-0{,}9M=0{,}3M=120\Rightarrow M=400\) TL. \ Kontrol: \(480-360=120\).
\end{cozum}

\begin{ornek}
Bir ürünün etiket fiyatı, maliyetinin \(\%60\) fazlasıdır. Ürün etiket fiyatı üzerinden \(\%25\) indirimle satılırsa maliyet üzerinden yüzde kaç kâr edilir?
\end{ornek}
\begin{cozum}
Maliyet \(100\) \(\Rightarrow\) etiket \(160\) \(\Rightarrow\) indirimli satış \(160\cdot 0{,}75=120\) \(\Rightarrow\) \(\%20\) kâr.
\end{cozum}

\begin{ornek}
Kilosu \(10\) TL'ye alınan yaş üzüm kuruyunca ağırlığının \(\%20\)'sini kaybediyor. Kuru üzümün kilosunun maliyeti kaç TL olur?
\end{ornek}
\begin{cozum}
\(100\) kg yaş üzüm \(1000\) TL'ye alınır. Kuruyunca \(80\) kg kalır; maliyet değişmez.
\[\frac{1000}{80}=12{,}5\ \text{TL/kg}\]
\end{cozum}

% =====================================================
\section{Karışım Problemleri}

\begin{kural}[Karışım ve tablo yöntemi]
\[\text{Oran}=\frac{\text{Madde miktarı}}{\text{Karışım miktarı}},\qquad \text{Madde miktarı}=\text{Karışım miktarı}\cdot\text{Oran}\]
Her kap için bir satır yazılır: \textit{karışım miktarı} \(\mid\) \textit{oran} \(\mid\) \textit{madde miktarı}. Kaplar birleştirilince karışım miktarları ve madde miktarları ayrı ayrı toplanır; yeni oran \(=\) toplam madde / toplam karışım (ağırlıklı ortalama).
\begin{itemize}[leftmargin=*,itemsep=1pt]
\item Saf su ekleme: madde sabit, karışım artar (suyun oranı \(\%0\)).
\item Su buharlaştırma: madde sabit, karışım azalır.
\item Saf madde ekleme: madde de karışım da artar (saf maddenin oranı \(\%100\)).
\end{itemize}
\end{kural}

\begin{dikkat}
Homojen bir karışımdan bir miktar \textbf{alınırsa (dökülürse) oran değişmez}; yalnız madde ve karışım miktarları aynı oranda azalır. Bu yüzden ``karışımdan kaç kg alınırsa oran değişir'' türü bir soru kendi başına anlamsızdır.
\end{dikkat}

\begin{ornek}
Tuz oranı \(\%25\) olan \(80\) kg tuzlu suya \(20\) kg saf su ekleniyor. Yeni tuz oranı yüzde kaçtır?
\end{ornek}
\begin{cozum}
Tuz: \(80\cdot 0{,}25=20\) kg. \ Yeni karışım \(100\) kg. \ \(\dfrac{20}{100}=\%20\).
\end{cozum}

\begin{ornek}
Şeker oranı \(\%15\) olan \(40\) kg şekerli sudan \(10\) kg su buharlaştırılıyor. Yeni şeker oranı yüzde kaçtır?
\end{ornek}
\begin{cozum}
Şeker: \(40\cdot 0{,}15=6\) kg (değişmez). Karışım \(30\) kg. \ \(\dfrac{6}{30}=\dfrac{20}{100}=\%20\).
\end{cozum}

\needspace{16\baselineskip}
\begin{ornek}
Tuz oranı \(\%20\) ve \(\%50\) olan iki tuzlu sudan, tuz oranı \(\%30\) olan \(90\) kg karışım hazırlanacaktır. Her birinden kaç kg kullanılmalıdır?
\end{ornek}
\begin{cozum}
\begin{center}
\begin{tabular}{lccc}
\toprule
 & Karışım (kg) & Oran & Tuz (kg)\\\midrule
1. kap & \(x\) & \(\%20\) & \(0{,}2x\)\\
2. kap & \(y\) & \(\%50\) & \(0{,}5y\)\\
Karışım & \(90\) & \(\%30\) & \(27\)\\\bottomrule
\end{tabular}
\end{center}
\(x+y=90\) ve \(0{,}2x+0{,}5y=27\). \ \(0{,}2(90-y)+0{,}5y=27\Rightarrow 18+0{,}3y=27\Rightarrow y=30\), \ \(x=60\) kg.
\end{cozum}

\begin{ornek}
Şeker oranı \(\%10\) olan \(50\) kg şekerli sudan \(10\) kg alınıyor; kalan karışıma \(10\) kg şeker ekleniyor. Son şeker oranı yüzde kaçtır?
\end{ornek}
\begin{cozum}
\(10\) kg alınınca oran değişmez: \(40\) kg karışımda \(4\) kg şeker. \(10\) kg şeker eklenince: \(\dfrac{4+10}{40+10}=\dfrac{14}{50}=\%28\).
\end{cozum}

\begin{kural}[Alaşım: altın ayarı]
Saf altın \(24\) ayardır. \(k\) ayar altında saf altın oranı \(\dfrac{k}{24}\)'tür: \quad saf altın \(=\) kütle \(\cdot\dfrac{k}{24}\). Alaşımlar karıştırılınca saf altın miktarları toplanır; yeni ayar \(=\dfrac{\text{toplam saf altın}}{\text{toplam kütle}}\cdot 24\).
\end{kural}

\begin{ornek}
\(14\) ayar \(30\) g altın ile \(22\) ayar \(18\) g altın eritilerek karıştırılıyor. Elde edilen altın kaç ayardır?
\end{ornek}
\begin{cozum}
Saf altın: \(30\cdot\dfrac{14}{24}=17{,}5\) g ve \(18\cdot\dfrac{22}{24}=16{,}5\) g; toplam \(34\) g. Toplam kütle \(48\) g.
\[\text{Ayar}=\frac{34}{48}\cdot 24=17\]
\end{cozum}

% =====================================================
\section{Faiz Problemleri}

\begin{kural}[Basit faiz]
\(A\): anapara, \(n\): \textbf{yıllık} faiz oranı (yüzde), \(F\): faiz.
\[F=\frac{A\cdot n\cdot t}{100}\ (t\ \text{yıl}),\qquad F=\frac{A\cdot n\cdot t}{1200}\ (t\ \text{ay}),\qquad F=\frac{A\cdot n\cdot t}{36000}\ (t\ \text{gün})\]
Gün hesabında \textbf{\(1\) yıl \(=360\) gün} kabul edilir (\(36000=100\cdot 360\)). Dönem sonunda çekilen toplam para \(A+F\)'dir.
\end{kural}

\begin{kural}[Bileşik faiz (kısa bilgi)]
Her yıl sonunda faiz anaparaya eklenir; sonraki yılın faizi bu toplam üzerinden hesaplanır. \(t\) yıl sonunda toplam para:
\(\ A\cdot\left(1+\dfrac{n}{100}\right)^t\). \ Soruda ``bileşik'' denmedikçe faiz basit faizdir.
\end{kural}

\begin{ornek}
\(15000\) TL yıllık \(\%32\) basit faizle \(9\) ay bankada kalırsa kaç TL faiz getirir? Dönem sonunda toplam kaç TL çekilir?
\end{ornek}
\begin{cozum}
\(F=\dfrac{15000\cdot 32\cdot 9}{1200}=3600\) TL. \quad Toplam: \(15000+3600=18600\) TL.
\end{cozum}

\begin{ornek}
\(7200\) TL yıllık \(\%30\) basit faizle \(75\) günde kaç TL faiz getirir? (\(1\) yıl \(=360\) gün)
\end{ornek}
\begin{cozum}
\(F=\dfrac{7200\cdot 30\cdot 75}{36000}=450\) TL.
\end{cozum}

\begin{ornek}
Yıllık \(\%25\) basit faizle bankaya yatırılan bir para kaç yıl sonra iki katına çıkar?
\end{ornek}
\begin{cozum}
İki katına çıkması için faiz anaparaya eşit olmalı: \(A=\dfrac{A\cdot 25\cdot t}{100}\Rightarrow 25t=100\Rightarrow t=4\) yıl.
\end{cozum}

\begin{ornek}
Basit faizle bankaya yatırılan bir para \(2\) yıl sonunda \(13000\) TL, \(5\) yıl sonunda \(17500\) TL olmaktadır. Anapara ve yıllık faiz oranı nedir?
\end{ornek}
\begin{cozum}
\(3\) yıllık faiz: \(17500-13000=4500\) TL \(\Rightarrow\) yıllık faiz \(1500\) TL. \ Anapara: \(13000-2\cdot 1500=10000\) TL.\\
\(\dfrac{10000\cdot n\cdot 1}{100}=1500\Rightarrow n=15\) \(\Rightarrow\) yıllık \(\%15\).
\end{cozum}

\begin{ornek}
\(20000\) TL yıllık \(\%10\) bileşik faizle \(3\) yıl bankada kalırsa toplam kaç TL olur? Basit faizle olsaydı ne kadar olurdu?
\end{ornek}
\begin{cozum}
\(20000\to 22000\to 24200\to 26620\) TL \ (\(20000\cdot 1{,}1^3=26620\)).\\
Basit faiz: \(F=\dfrac{20000\cdot 10\cdot 3}{100}=6000\) \(\Rightarrow\) \(26000\) TL. Fark: \(620\) TL.
\end{cozum}

% =====================================================
\section{Özet: Formül Kartı}
\begin{center}
\renewcommand{\arraystretch}{1.6}
\begin{tabular}{@{}ll@{}}
\toprule
Konu & Formül / Kural\\\midrule
Yüzde & \(x\)'in \(\%a\)'sı \(=x\cdot\frac{a}{100}\); \ \(A\), \(B\)'nin \(\frac{A}{B}\cdot 100\) yüzdesi\\
Artış / azalış & \(x\left(1+\frac{a}{100}\right)\) \ / \ \(x\left(1-\frac{a}{100}\right)\); ardışık değişimde çarpanlar çarpılır\\
Kâr -- zarar & Satış \(=\) Maliyet \(\pm\) Kâr/Zarar; \ yüzde \textbf{maliyet} üzerinden\\
İndirim & İndirimli fiyat \(=\) Etiket \(\cdot\left(1-\frac{d}{100}\right)\); \ yüzde \textbf{etiket} üzerinden\\
Karışım & Madde \(=\) Karışım \(\cdot\) Oran; \ yeni oran \(=\) toplam madde / toplam karışım\\
Su / buhar / saf madde & su ve buharda madde sabit; saf maddenin oranı \(\%100\)\\
Altın ayarı & saf altın \(=\) kütle \(\cdot\frac{\text{ayar}}{24}\)\\
Basit faiz & \(F=\frac{A\,n\,t}{100}\) (yıl), \(\frac{A\,n\,t}{1200}\) (ay), \(\frac{A\,n\,t}{36000}\) (gün; yıl \(=360\) gün)\\
Bileşik faiz & Toplam \(=A\left(1+\frac{n}{100}\right)^t\)\\\bottomrule
\end{tabular}
\end{center}

% =====================================================
\newpage
% Hafta 14 -- alıştırma gövdesi (ders notu ve ayrı alıştırma kağıdı tarafından \input edilir)
\section{Çözümlü Alıştırmalar}
Her soruyu önce kendiniz deneyin, sonra çözümü okuyun.

\begin{alistirma}
Bir sayının \(\%25\)'inin \(\%40\)'ı \(18\) ise bu sayı kaçtır?
\end{alistirma}
\begin{cozum}
Sayı \(x\) olsun. \(\%25\)'i \(0{,}25x\); bunun \(\%40\)'ı \(0{,}40\cdot 0{,}25x=0{,}10x\).
\[0{,}10x=18\ \Rightarrow\ x=180\]
Kontrol: \(180\)'in \(\%25\)'i \(45\), \(45\)'in \(\%40\)'ı \(18\) \(\checkmark\).
\end{cozum}

\begin{alistirma}
Bir ürünün fiyatına önce \(\%25\) zam, ardından yeni fiyat üzerinden \(\%20\) zam yapılıyor. Son fiyat ilk fiyata göre yüzde kaç artmıştır?
\end{alistirma}
\begin{cozum}
İlk fiyat \(100\) olsun: \(100\xrightarrow{+\%25}125\xrightarrow{+\%20\ (25)}150\).\\
Son fiyat \(150\) \(\Rightarrow\) \(\%50\) artış. (\(\%25+\%20=\%45\) değil; çarpanlar: \(1{,}25\cdot 1{,}2=1{,}5\).)
\end{cozum}

\begin{alistirma}
Bir sınıftaki öğrencilerin \(\%60\)'ı kızdır. Kızların \(\%30\)'u, erkeklerin \(\%45\)'i basketbol oynamaktadır. Sınıfın yüzde kaçı basketbol oynamaktadır?
\end{alistirma}
\begin{cozum}
Sınıf \(100\) kişi olsun: \(60\) kız, \(40\) erkek.
\[60\cdot\frac{30}{100}+40\cdot\frac{45}{100}=18+18=36\ \Rightarrow\ \%36\]
\end{cozum}

\begin{alistirma}
\(\%25\) kârla \(400\) TL'ye satılan bir ürün \(\%10\) zararla satılsaydı kaç TL'ye satılırdı?
\end{alistirma}
\begin{cozum}
Maliyet \(100x\) \(\Rightarrow\) \(\%25\) kârla satış \(125x=400\Rightarrow x=3{,}2\). Maliyet \(=320\) TL.\\
\(\%10\) zararla satış: \(320\cdot 0{,}90=288\) TL.
\end{cozum}

\begin{alistirma}
Bir mağaza etiket fiyatı üzerinden \(\%30\) indirim yaptığında bile maliyet üzerinden \(\%12\) kâr etmektedir. Mağaza indirimden önce maliyet üzerinden yüzde kaç kârla satış yapıyordu?
\end{alistirma}
\begin{cozum}
Maliyet \(100\) \(\Rightarrow\) indirimli satış \(112\). Etiket fiyatı \(E\) ise
\[E\cdot 0{,}70=112\ \Rightarrow\ E=160\]
Etiket fiyatı maliyetin \(\%160\)'ı \(\Rightarrow\) indirimden önce \(\%60\) kâr.
\end{cozum}

\begin{alistirma}
Alkol oranı \(\%40\) olan \(30\) L kolonyaya kaç L saf su eklenirse alkol oranı \(\%25\) olur?
\end{alistirma}
\begin{cozum}
Alkol miktarı değişmez: \(30\cdot 0{,}40=12\) L. Eklenen su \(x\) L:
\[\frac{12}{30+x}=\frac{25}{100}\ \Rightarrow\ 30+x=48\ \Rightarrow\ x=18\ \text{L}\]
\end{cozum}

\begin{alistirma}
Tuz oranı \(\%20\) olan \(40\) kg tuzlu su ile tuz oranı \(\%30\) olan \(60\) kg tuzlu su karıştırılıyor. Karışımdan \(35\) kg su buharlaştırılırsa son tuz oranı yüzde kaç olur?
\end{alistirma}
\begin{cozum}
Tuz: \(40\cdot 0{,}20+60\cdot 0{,}30=8+18=26\) kg; karışım \(100\) kg.\\
Buharlaştırmada tuz değişmez; karışım \(100-35=65\) kg.
\[\frac{26}{65}=\frac{2}{5}=\frac{40}{100}\ \Rightarrow\ \%40\]
\end{cozum}

\begin{alistirma}
\(8000\) TL'nin bir kısmı yıllık \(\%30\), kalanı yıllık \(\%20\) basit faizle bankaya yatırılıyor. \(1\) yıl sonunda toplam \(2000\) TL faiz alındığına göre \(\%30\) faizle kaç TL yatırılmıştır?
\end{alistirma}
\begin{cozum}
\(\%30\)'a yatırılan \(x\) TL, \(\%20\)'ye yatırılan \(8000-x\) TL:
\[\frac{30x}{100}+\frac{20(8000-x)}{100}=2000\ \Rightarrow\ 30x+160000-20x=200000\ \Rightarrow\ x=4000\ \text{TL}\]
Kontrol: \(1200+800=2000\) \(\checkmark\).
\end{cozum}

\kendinizoncesi
\section{Alıştırmalar (Kendiniz Çözün)}

\begin{enumerate}[label=\textbf{\arabic*.},leftmargin=*,itemsep=6pt]
\item \(640\) sayısının \(\%15\)'i kaçtır?
\item \(36\) sayısı \(240\)'ın yüzde kaçıdır?
\item Bir sayının \(\%35\)'i \(84\) ise bu sayı kaçtır?
\item Fiyatı \(1200\) TL olan bir ürüne önce \(\%15\) zam, sonra yeni fiyat üzerinden \(\%10\) indirim yapılıyor. Son fiyat kaç TL'dir?
\item \(\%30\) indirimle \(560\) TL'ye satılan bir ürünün indirimsiz fiyatı kaç TL'dir?
\item Bir dikdörtgenin uzun kenarı \(\%25\) artırılıyor. Alanın değişmemesi için kısa kenar yüzde kaç azaltılmalıdır?
\item Maliyeti \(250\) TL olan bir ürün \(310\) TL'ye satılıyor. Kâr yüzdesi kaçtır?
\item \(5\) tanesi \(20\) TL'ye alınan yumurtaların \(4\) tanesi \(20\) TL'ye satılıyor. Yüzde kaç kâr edilir?
\item Etiket fiyatı \(900\) TL olan bir ürün \(\%20\) indirimle satıldığında maliyet üzerinden \(\%20\) kâr ediliyor. Ürünün maliyeti kaç TL'dir?
\item Şeker oranı \(\%40\) olan \(25\) kg şekerli suya kaç kg su eklenirse şeker oranı \(\%25\) olur?
\item Tuz oranı \(\%20\) olan \(60\) kg tuzlu sudan kaç kg su buharlaştırılırsa tuz oranı \(\%30\) olur?
\item Alkol oranı \(\%15\) olan \(20\) L ile \(\%35\) olan \(30\) L kolonya karıştırılıyor. Yeni alkol oranı yüzde kaçtır?
\item \(22\) ayar \(36\) g altına kaç g saf (\(24\) ayar) altın eklenirse \(23\) ayar altın elde edilir?
\item \(4500\) TL yıllık \(\%40\) basit faizle \(8\) ayda kaç TL faiz getirir?
\item \(6000\) TL basit faizle \(3\) yıl sonunda \(9600\) TL'ye ulaşıyor. Yıllık faiz oranı yüzde kaçtır?
\item \(5000\) TL yıllık \(\%10\) bileşik faizle \(2\) yıl sonunda kaç TL olur?
\end{enumerate}

\subsection*{Cevap Anahtarı}
{\small
\begin{multicols}{2}
\begin{enumerate}[label=\textbf{\arabic*.},leftmargin=*,itemsep=1pt]
\item \(640\cdot 0{,}15=96\).
\item \(\frac{36}{240}\cdot 100=\%15\).
\item \(0{,}35x=84\Rightarrow x=240\).
\item \(1200\cdot 1{,}15\cdot 0{,}9=1242\) TL.
\item \(0{,}7x=560\Rightarrow x=800\) TL.
\item \(1{,}25\cdot k=1\Rightarrow k=0{,}8\): \(\%20\) azaltılmalı.
\item Kâr \(60\); \(\frac{60}{250}\cdot 100=\%24\).
\item \(20\) yumurta: maliyet \(80\), satış \(100\) \(\Rightarrow\) \(\%25\).
\item Satış \(720\); \(1{,}2M=720\Rightarrow M=600\) TL.
\item Şeker \(10\) kg; \(\frac{10}{25+x}=0{,}25\Rightarrow x=15\) kg.
\item Tuz \(12\) kg; \(\frac{12}{60-x}=0{,}3\Rightarrow x=20\) kg.
\item \(\frac{3+10{,}5}{50}=\frac{27}{100}\Rightarrow\%27\).
\item Saf: \(33+x=\frac{23}{24}(36+x)\Rightarrow x=36\) g.
\item \(\frac{4500\cdot 40\cdot 8}{1200}=1200\) TL.
\item \(F=3600=\frac{6000\cdot n\cdot 3}{100}\Rightarrow n=20\): \(\%20\).
\item \(5000\cdot 1{,}1^2=6050\) TL.
\end{enumerate}
\end{multicols}}

% =====================================================
\newcounter{karma}
\newenvironment{karma}[1]{\refstepcounter{karma}\par\medskip\needspace{9\baselineskip}\noindent\textbf{\color{anamavi}Soru \thekarma.}\ {\small\textit{(#1)}}\ }{\par}
\needspace{15\baselineskip}
\section{Finale Hazırlık: Karma Tekrar (Hafta 1--14)}
Bu bölüm final sınavı kapsamındaki tüm konulardan sorular içerir. İlk \(6\) soru tam çözümlüdür; kalan \(10\) soru için bölüm sonunda cevap anahtarı vardır.

\subsection*{Tam Çözümlü Sorular}

\begin{karma}{Hafta 1: faktöriyel}
\(\dfrac{8!-7!}{6!}\) işleminin sonucu kaçtır?
\end{karma}
\begin{cozum}
\(8!=8\cdot 7!\) olduğundan \(8!-7!=7!\,(8-1)=7\cdot 7!\). \ \(7!=7\cdot 6!\):
\[\frac{7\cdot 7\cdot 6!}{6!}=49\]
\end{cozum}

\begin{karma}{Hafta 3: bölünebilme}
\(a\) ve \(b\) rakam olmak üzere, dört basamaklı \(\ab{4a7b}\) sayısı \(45\) ile tam bölünebilmektedir. \(a\)'nın alabileceği değerlerin toplamı kaçtır?
\end{karma}
\begin{cozum}
\(45=5\cdot 9\) ve \(5\) ile \(9\) aralarında asal: sayı hem \(5\)'e hem \(9\)'a bölünmeli.\\
\(5\)'e bölünebilme: \(b=0\) veya \(b=5\). \quad \(9\)'a bölünebilme: \(4+a+7+b=11+a+b\), \(9\)'un katı.\\
\(b=0\): \(11\le 11+a\le 20\) aralığında \(9\)'un tek katı \(18\) \(\Rightarrow a=7\) \ (\(4770=45\cdot 106\)).\\
\(b=5\): \(16\le 16+a\le 25\) aralığında \(9\)'un tek katı \(18\) \(\Rightarrow a=2\) \ (\(4275=45\cdot 95\)).\\
\(a\)'nın değerleri \(7\) ve \(2\) \(\Rightarrow\) toplam \(9\).
\end{cozum}

\begin{karma}{Hafta 7: mutlak değer}
\(|x-2|<5\) eşitsizliğini sağlayan kaç tam sayı vardır? Bu tam sayıların toplamı kaçtır?
\end{karma}
\begin{cozum}
\(|x-2|<5\iff -5<x-2<5\iff -3<x<7\).\\
Tam sayılar: \(-2,-1,0,1,2,3,4,5,6\) \(\Rightarrow\) \(9\) tane. \ Toplam: \((-2)+(-1)+0+1+2+\dots+6=18\).
\end{cozum}

\begin{karma}{Hafta 11: çarpanlara ayırma}
\(x\neq 2\) ve \(x\neq -3\) olmak üzere \(\dfrac{x^2-9}{x^2+x-6}\) ifadesinin en sade biçimi nedir?
\end{karma}
\begin{cozum}
Pay: iki kare farkı, \(x^2-9=(x-3)(x+3)\). \ Payda: çarpımı \(-6\), toplamı \(1\) olan sayılar \(3\) ve \(-2\): \(x^2+x-6=(x+3)(x-2)\).
\[\frac{(x-3)(x+3)}{(x+3)(x-2)}=\frac{x-3}{x-2}\]
\end{cozum}

\begin{karma}{Hafta 14: kâr -- zarar ve indirim}
Bir mağaza bir ürünün etiket fiyatını maliyetin \(\%50\) fazlası olarak belirliyor. Ürün etiket fiyatı üzerinden \(\%20\) indirimle \(360\) TL'ye satılıyor. Ürünün maliyeti kaç TL'dir ve bu satıştan yüzde kaç kâr edilmiştir?
\end{karma}
\begin{cozum}
Etiket \(E\): \ \(E\cdot 0{,}80=360\Rightarrow E=450\) TL. \quad Maliyet \(M\): \ \(1{,}5M=450\Rightarrow M=300\) TL.\\
Kâr: \(360-300=60\) TL \(\Rightarrow\) \(\dfrac{60}{300}\cdot 100=\%20\) kâr.
\end{cozum}

\begin{karma}{Hafta 13: yaş problemi}
Bir annenin bugünkü yaşı, iki çocuğunun yaşları toplamının \(3\) katıdır. \(6\) yıl sonra annenin yaşı, çocuklarının o günkü yaşları toplamının \(2\) katı olacaktır. Anne bugün kaç yaşındadır?
\end{karma}
\begin{cozum}
Çocukların yaşları toplamı bugün \(s\), annenin yaşı \(3s\). \(6\) yıl sonra anne \(3s+6\); iki çocuğun her biri \(6\) yaş büyüyeceğinden toplam \(s+12\).
\[3s+6=2(s+12)\ \Rightarrow\ 3s+6=2s+24\ \Rightarrow\ s=18\]
Anne bugün \(3\cdot 18=54\) yaşındadır. Kontrol: \(6\) yıl sonra \(60=2\cdot 30\) \(\checkmark\).
\end{cozum}

\subsection*{Cevap Anahtarlı Sorular}

\begin{enumerate}[label=\textbf{\arabic*.},leftmargin=*,itemsep=6pt]\setcounter{enumi}{6}
\item \textit{(Hafta 2)} \(180\) sayısının kaç pozitif böleni vardır?
\item \textit{(Hafta 3)} Kenarları \(48\) m ve \(72\) m olan dikdörtgen bir bahçe, eşit büyüklükte kare parsellere ayrılacaktır. En az kaç parsel elde edilir?
\item \textit{(Hafta 4)} \(0{,}5+\dfrac{2}{3}\cdot\dfrac{3}{4}\) işleminin sonucu kaçtır?
\item \textit{(Hafta 5)} \(3(x-2)+4=2x+5\) denkleminin çözümü nedir?
\item \textit{(Hafta 6)} Bir otoparkta otomobil ve motosikletlerden oluşan \(40\) araç vardır. Tekerleklerin toplam sayısı \(130\) olduğuna göre kaç motosiklet vardır? (Otomobil \(4\), motosiklet \(2\) tekerleklidir.)
\item \textit{(Hafta 9)} \(9^{x}=27^{x-1}\) ise \(x\) kaçtır?
\item \textit{(Hafta 10)} \(\dfrac{6}{\sqrt{3}}+\sqrt{27}\) işleminin sonucu kaçtır?
\item \textit{(Hafta 12)} \(5\) sayının aritmetik ortalaması \(18\)'dir. Bu sayılardan biri çıkarılınca kalan \(4\) sayının ortalaması \(16\) oluyor. Çıkarılan sayı kaçtır?
\item \textit{(Hafta 14)} Tuz oranı \(\%20\) olan \(30\) kg tuzlu suya \(10\) kg tuz ekleniyor. Yeni tuz oranı yüzde kaçtır?
\item \textit{(Hafta 14)} \(7200\) TL yıllık \(\%25\) basit faizle \(120\) günde kaç TL faiz getirir? (\(1\) yıl \(=360\) gün)
\end{enumerate}

\subsection*{Cevap Anahtarı (Karma Tekrar)}
{\small
\begin{multicols}{2}
\begin{enumerate}[label=\textbf{\arabic*.},leftmargin=*,itemsep=1pt]
\item \(49\) \ (tam çözümlü)
\item \(9\) \ (tam çözümlü)
\item \(9\) tam sayı; toplam \(18\) \ (tam çözümlü)
\item \(\frac{x-3}{x-2}\) \ (tam çözümlü)
\item Maliyet \(300\) TL; \(\%20\) kâr \ (tam çözümlü)
\item \(54\) \ (tam çözümlü)
\item \(180=2^2\cdot 3^2\cdot 5\Rightarrow 3\cdot 3\cdot 2=18\).
\item EBOB\((48,72)=24\); \(\frac{48}{24}\cdot\frac{72}{24}=2\cdot 3=6\) parsel.
\item \(0{,}5+0{,}5=1\).
\item \(3x-2=2x+5\Rightarrow x=7\).
\item \(4a+2m=130,\ a+m=40\Rightarrow a=25,\ m=15\).
\item \(3^{2x}=3^{3x-3}\Rightarrow 2x=3x-3\Rightarrow x=3\).
\item \(2\sqrt{3}+3\sqrt{3}=5\sqrt{3}\).
\item \(5\cdot 18-4\cdot 16=90-64=26\).
\item \(\frac{6+10}{30+10}=\frac{16}{40}\Rightarrow\%40\).
\item \(\frac{7200\cdot 25\cdot 120}{36000}=600\) TL.
\end{enumerate}
\end{multicols}}


\vspace{6pt}
\noindent\textcolor{anamavi}{\rule{\textwidth}{0.5pt}}\\
{\small Gelecek hafta: FİNAL SINAVI (Hafta 1--14 konuları).}

\end{document}
